\documentclass[11pt,a4paper]{article}
\usepackage[T1]{fontenc}
\usepackage[utf8]{inputenc}
\usepackage{amsmath,amssymb,amsthm}
\usepackage[margin=27mm]{geometry}
\usepackage[hidelinks]{hyperref}
\newtheorem{proposition}{Proposizione}[section]
\newtheorem{theorem}[proposition]{Teorema}
\newtheorem{lemma}[proposition]{Lemma}
\newtheorem{corollary}[proposition]{Corollario}
\theoremstyle{remark}\newtheorem{remark}[proposition]{Osservazione}
\newcommand{\E}{\mathbb E_p}
\newcommand{\Pp}{\mathbb P_p}
\newcommand{\ind}{\mathbf1}
\newcommand{\Mmax}{M_{\max}}
\newcommand{\F}{\mathcal F}
\title{Lower confidence bounds per la massima prevalenza non osservata}
\author{Note teoriche: rapporti di momenti e controllo simultaneo}
\date{6 novembre 2026}
\begin{document}
\maketitle
\tableofcontents

\section{Modello, notazione e fatti probabilistici comuni}
Il tempo di campionamento $n$ è fissato. Osserviamo variabili indipendenti
$N_j\sim\mathrm{Bin}(n,p_j)$, $j\ge1$, con parametro fissato
\[
 p\in\F:=\{(p_j)_{j\ge1}:0\le p_j\le1,\quad
 S(p):=\sum_{j\ge1}p_j<\infty\}.
\]
Definiamo, per $a\ge1$ e $1\le k\le n$,
\[
 \begin{gathered}
 M_a=\sum_jp_j^a\ind\{N_j=0\},\qquad \Mmax=\max_{j:N_j=0}p_j,\\
 K_k=K_{n,k}=\sum_j\ind\{N_j=k\},\qquad \eta_k=\E K_k.
 \end{gathered}
\]
Il massimo è zero se non vi sono prevalenze positive non osservate;
la missing mass è $\mathcal M_0=M_1$. Le statistiche $M_a$ e $\Mmax$
sono casuali, benché $p$ sia fissato. Tutte le garanzie sono a tempo $n$
fissato e non sono garanzie di arresto opzionale.
Poiché $\E\sum_jN_j=nS(p)$, i conteggi osservati e gli $\eta_k$ sono finiti.
Se $\Mmax>0$, il massimo è raggiunto e ha molteplicità finita:
la sommabilità implica che vi siano solo finitamente molti $p_j$ sopra
qualsiasi soglia positiva.

Per interi $1\le j\le n$, poniamo
\[
 a_j=\binom nj^{-1},\quad T_j=a_jK_j,\quad D_j=T_j-M_j,
 \quad b_j=\E D_j,\quad e_j=D_j-b_j.
\]
\begin{proposition}[Bias e varianza dell'errore]\label{prop:errors}
Se $n\ge2j$, allora
\[
 0\le b_j\le \overline b_j:=j a_{j+1}\eta_{j+1},\qquad
 \operatorname{Var}_p(D_j)\le v_j:=a_j^2\eta_j+a_{2j}\eta_{2j}.
\]
Per ogni $x>0$ ed $\epsilon\in(0,1)$,
\begin{align}
 \Pp\{D_j>\overline b_j+\sqrt{2v_jx}+2a_jx/3\}&\le e^{-x},\label{eq:bern}\\
 \Pp\{D_j<-\sqrt{(1-\epsilon)v_j/\epsilon}\}&\le\epsilon.\label{eq:cant}
\end{align}
\end{proposition}
\begin{proof}
Scriviamo $q=1-p$ per una singola categoria e
$d_j=a_j\ind\{N=j\}-p^j\ind\{N=0\}$. Allora
\[
 \mathbb E d_j=p^jq^{n-j}(1-q^j)
 \le j p^{j+1}q^{n-j-1},
\]
poiché $1-q^j\le jp$ e $q^{n-j}\le q^{n-j-1}$.
Sommando otteniamo il controllo del bias. Per disgiunzione degli eventi,
\[
 \sum_\ell\mathbb E d_{j,\ell}^2
 =a_j^2\eta_j+\sum_\ell p_\ell^{2j}(1-p_\ell)^n
 \le a_j^2\eta_j+a_{2j}\eta_{2j}.
\]
L'indipendenza dà il controllo della varianza. Inoltre
$d_{j,\ell}-\mathbb E d_{j,\ell}\le a_j$, poiché la media è non negativa.
Bernstein unilaterale dà \eqref{eq:bern}. Per \eqref{eq:cant},
$b_j\ge0$ e Cantelli implicano
$\Pp(D_j<-t)\le v_j/(v_j+t^2)$.
Se $v_j=0$, $D_j=b_j\ge0$ quasi certamente e la conclusione è diretta.
Le somme infinite si ottengono per passaggio al limite: la somma dei valori
attesi assoluti è finita e la somma delle varianze è finita.
\end{proof}

\begin{lemma}[Upper bound osservabile dei conteggi medi]\label{lem:counts}
Per $k\ge1$ e $\epsilon\in(0,1)$, definiamo
\[
 C_k(\epsilon)=\left(\sqrt{\frac{\log(1/\epsilon)}2}
 +\sqrt{K_k+\frac76\log(1/\epsilon)}\right)^2.
\]
Allora $\Pp\{\eta_k\le C_k(\epsilon)\}\ge1-\epsilon$.
Per un insieme deterministico finito $J$ di $m$ indici,
con probabilità almeno $1-m\epsilon$ tutti questi bound valgono insieme.
\end{lemma}
\begin{proof}
$K_k$ è somma di Bernoulli indipendenti e ha varianza al più $\eta_k$.
Bernstein implica, con $x=\log(1/\epsilon)$,
$\Pp\{\eta_k-K_k>\sqrt{2\eta_kx}+2x/3\}\le\epsilon$.
Risolvendo la corrispondente disuguaglianza quadratica in $\sqrt{\eta_k}$
si ottiene $C_k$. L'ultima affermazione segue dall'unione degli eventi.
È il controllo dei conteggi usato anche da Ayed et al. (2019, Remark A.1).
\end{proof}

\section{Geometria oracle e scelta dei rapporti consecutivi}
In questa sezione fissiamo anche il campione e supponiamo $M_1>0$.
\begin{proposition}[Gerarchia oracle]\label{prop:oracle}
Per reali $1\le s<r$, poniamo
$Q_{r,s}=(M_r/M_s)^{1/(r-s)}$. Allora $Q_{r,s}\le\Mmax$.
La quantità $Q_{r,s}$ è non decrescente in ciascun estremo, mantenendo $s<r$.
Per interi $r\ge2$,
\[
 Q_{r,1}\le Q_{r,2}\le\cdots\le Q_{r,r-1}
 =\frac{M_r}{M_{r-1}}\le\Mmax,
 \qquad \frac{M_r}{M_{r-1}}\uparrow\Mmax.
\]
Tra le coppie intere $r>s\ge1$, il minimo oracle è $M_2/M_1$.
\end{proposition}
\begin{proof}
Si ha $M_r\le\Mmax^{r-s}M_s$. Per Hölder $f(t)=\log M_t$ è convessa;
$\log Q_{r,s}$ è la pendenza della secante tra $s$ e $r$ e cresce in
ciascun estremo. Equivalentemente, posto $A_j=M_j/M_{j-1}$,
la log-convessità dà $A_j\le A_{j+1}$ e
$Q_{r,s}=(\prod_{j=s+1}^r A_j)^{1/(r-s)}$.
Per il limite, sia $m=\Mmax>0$ di molteplicità $k$. Scriviamo
$M_j=m^j(k+\sum_{0<u_\ell<m}(u_\ell/m)^j)$, dove gli $u_\ell$
sono le prevalenze non osservate. La somma residua tende a zero per
convergenza dominata, essendo dominata da $\sum u_\ell/m<\infty$.
Ne segue $A_j\to m$. Se tutte le prevalenze positive non osservate
sono uguali, tutte le disuguaglianze sono uguaglianze.
\end{proof}
\begin{proposition}[Media pesata e quoziente di Rayleigh]
Con $u_j=p_j\ind\{N_j=0\}$, per $r\ge2$,
\[
 \frac{M_r}{M_{r-1}}=\sum_j\frac{u_j^{r-1}}{M_{r-1}}u_j
 =\frac{\langle v,\operatorname{diag}(u)v\rangle}{\langle v,v\rangle},
 \qquad v_j=u_j^{(r-1)/2}.
\]
\end{proposition}
\begin{proof}
Le identità sono immediate. La sommabilità assicura $v\in\ell_2$.
La norma di questo operatore è il massimo degli $u_j$, cioè $\Mmax$.
\end{proof}
La gerarchia riguarda i momenti veri. Non implica che i rapporti stimati,
dopo correzioni di bias e concentrazione, siano ordinati nello stesso modo.

\section{Prima derivazione: rapporto di estremi separati, ordine fissato}
\begin{theorem}[Lower bound separato]\label{thm:separate}
Fissiamo un intero $r\ge2$, $s=r-1$, $n\ge2r$ e $\delta\in(0,1)$.
Siano
\[
 J=\{s,r,r+1,2s,2r\},\quad m=|J|,\quad
 \epsilon=\frac{\delta}{m+2},\quad x=\log(1/\epsilon),\quad
 C_k=C_k(\epsilon)\ (k\in J).
\]
Per $j=s,r$ definiamo
\[
 B_j=j a_{j+1}C_{j+1},\qquad V_j=a_j^2C_j+a_{2j}C_{2j}.
\]
Tutti gli indici richiesti sono in $J$. Definiamo
\[
 L_r=[T_r-B_r-\sqrt{2V_rx}-2a_rx/3]_+,\qquad
 U_s=T_s+\sqrt{(1-\epsilon)V_s/\epsilon}.
\]
Allora $\Pp\{L_r\le M_r,\ M_s\le U_s\}\ge1-\delta$ e
\[
 \Pp\{\Mmax\ge L_r/U_s\}\ge1-\delta.
\]
Il rapporto è definito zero se $U_s=0$.
\end{theorem}
\begin{proof}
Il Lemma \ref{lem:counts} controlla $m$ eventi. Aggiungiamo
\eqref{eq:bern} per $D_r$ e \eqref{eq:cant} per $D_s$, ciascuno con
errore $\epsilon$. Sul loro evento congiunto $b_j\le B_j$, $v_j\le V_j$,
e le formule danno gli estremi dichiarati. Se $M_s>0$,
$\Mmax\ge M_r/M_s\ge L_r/U_s$. Se $M_s=0$, anche $M_r=0$
e sullo stesso evento $L_r=0$. L'errore totale è $(m+2)\epsilon$.
\end{proof}
\begin{remark}
Questo teorema conserva la derivazione con Cantelli per il denominatore.
Il benchmark già implementato $L_2/U_1$ usa invece, per $U_1$, lo
specifico upper bound della missing mass di Ayed et al. (2019).
Le due procedure non sono identiche e vengono tenute distinte.
\end{remark}

\section{Seconda derivazione: regione congiunta per una coppia fissata}
\begin{theorem}[Regione congiunta e controllo uniforme in $t$]\label{thm:pair}
Fissiamo $r\ge2$, $s=r-1$, $n\ge2r$, $\delta\in(0,1)$ e lo stesso
insieme $J$ del Teorema \ref{thm:separate}. Ora poniamo
\[
 m=|J|,\quad C_k=C_k(\delta/(2m)),\quad
 B_j=j a_{j+1}C_{j+1},\quad V_j=a_j^2C_j+a_{2j}C_{2j},
 \quad c=2/\sqrt\delta.
\]
Con probabilità almeno $1-\delta$, simultaneamente
\[
 0\le b_j\le B_j\ (j=s,r),\qquad
 \frac{e_s^2}{V_s}+\frac{e_r^2}{V_r}\le c^2.
\]
In particolare $(M_s,M_r)$ appartiene alla regione osservabile
\[
 \mathcal C_r=\left\{(u,v):\begin{gathered}
 0\le v\le u,\quad \exists\beta_s\in[0,B_s],\ \beta_r\in[0,B_r]:\\
 \frac{(T_s-u-\beta_s)^2}{V_s}+
 \frac{(T_r-v-\beta_r)^2}{V_r}\le c^2
 \end{gathered}\right\}.
\]
Sullo stesso evento, per ogni $t\ge0$,
\[
 M_r-tM_s\ge F_r(t):=T_r-tT_s-B_r-c\sqrt{V_r+t^2V_s}.
\]
Definito
\[
 \tau_r=\sup\{t\in[0,1]:F_r(t)>0\},\qquad \sup\varnothing:=0,
\]
si ha $\Pp\{\tau_r\le\Mmax\}\ge1-\delta$.
\end{theorem}
\begin{proof}
Se $v_j=0$, poniamo $e_j^2/v_j=0$: infatti $e_j=0$ quasi certamente.
La Proposizione \ref{prop:errors} dà
$\E(e_s^2/v_s+e_r^2/v_r)\le2$. Markov controlla la somma con soglia
$4/\delta$ ed errore $\delta/2$. Sull'evento dei conteggi, di errore
al più $\delta/2$, $V_j\ge v_j$ e $B_j\ge b_j$. Per unione dei due
eventi segue la regione. Non si è condizionato sui $V_j$ casuali.
Per Cauchy--Schwarz, contemporaneamente per tutti i $t$,
$e_r-te_s\le c\sqrt{V_r+t^2V_s}$. Inoltre $b_r-tb_s\le B_r$.
Se $F_r(t)>0$, allora $M_r-tM_s>0$: ciò implica $M_s>0$ e
$\Mmax\ge M_r/M_s>t$. Prendendo il sup si ottiene il risultato.
Questo argomento tratta anche $M_s=0$: in quel caso nessun $t$ può
essere certificato sul buon evento.
\end{proof}

\begin{proposition}[Soglia esplicita e confronto geometrico]\label{prop:root}
Per $V_r,V_s>0$, poniamo $A=T_r-B_r$, $D=T_s$, $v=c^2V_r$,
$w=c^2V_s$. Se $A\le\sqrt v$, allora $\tau_r=0$. Altrimenti,
\[
 N=A^2-v>0,\qquad
 \tau_r=\min\left\{1,\frac{N}{AD+\sqrt{D^2v+wN}}\right\}.
\]
Inoltre $\tau_r$ è almeno il bound rettangolare ottenuto dalla stessa
ellisse e con lo stesso $c$:
\[
 \tau_r\ge\min\left\{1,\frac{[T_r-B_r-c\sqrt{V_r}]_+}{T_s+c\sqrt{V_s}}\right\}.
\]
\end{proposition}
\begin{proof}
$F_r(t)=A-Dt-\sqrt{v+wt^2}$ è continua e strettamente decrescente
su $[0,\infty)$ (salvo derivata nulla in zero quando $D=0$).
Se $F_r(0)>0$ esiste una radice positiva unica. Elevando al quadrato
$A-Dt=\sqrt{v+wt^2}$ e razionalizzando la radice si ottiene la formula;
essa ha $A-Dt\ge0$ e quindi risolve anche l'equazione non quadratica.
Il confronto segue da $\sqrt{V_r+t^2V_s}\le\sqrt{V_r}+t\sqrt{V_s}$,
troncando entrambe le soglie a uno.
\end{proof}
\begin{remark}
Più precisamente, il confronto vale con il minimo tra uno e il rapporto
rettangolare. Nel confronto tra procedure non troncate su $[0,1]$,
si usano le radici non troncate. La regione è diagonale: non sfrutta
esplicitamente la covarianza. Markov introduce penalità proporzionali
a $\delta^{-1/2}$; il risultato non garantisce un miglioramento rispetto
al diverso benchmark Bernstein--Ayed.
\end{remark}

\section{Terza derivazione: controllo simultaneo degli ordini da 2 a R}
\begin{theorem}[Massimo adattivo con copertura simultanea]\label{thm:simultaneous}
Siano $R\ge2$ un intero deterministico, $n\ge2R$, $\delta\in(0,1)$.
Per $k=1,\ldots,2R$ e $j=1,\ldots,R$ definiamo
\[
 C_k=C_k(\delta/(4R)),\quad
 T_j=a_jK_j,\quad B_j=j a_{j+1}C_{j+1},\quad
 V_j=a_j^2C_j+a_{2j}C_{2j},\quad c_R=\sqrt{2R/\delta}.
\]
Per ogni $r=2,\ldots,R$ poniamo
\[
 F_{r,R}(t)=T_r-tT_{r-1}-B_r-c_R\sqrt{V_r+t^2V_{r-1}},
 \qquad \tau_{r,R}=\sup\{t\in[0,1]:F_{r,R}(t)>0\},
\]
con sup dell'insieme vuoto uguale a zero. Allora
\[
 \boxed{\inf_{p\in\F}\Pp\left\{\Mmax\ge
 \Lambda_R:=\max_{2\le r\le R}\tau_{r,R}\right\}\ge1-\delta.}
\]
La formula della Proposizione \ref{prop:root}, con $c=c_R$ e i $B,V$
qui definiti, calcola ciascuna soglia senza ottimizzazione numerica in $t$.
\end{theorem}
\begin{proof}
Il Lemma \ref{lem:counts} controlla tutti i $2R$ conteggi con errore
totale $\delta/2$. Indipendentemente da questo evento, nel senso che
non occorre condizionare, si ha
\[
 \E\sum_{j=1}^R\frac{e_j^2}{v_j}\le R,\qquad
 \Pp\left\{\sum_{j=1}^R\frac{e_j^2}{v_j}>\frac{2R}{\delta}\right\}
 \le\delta/2.
\]
Sull'intersezione, di probabilità almeno $1-\delta$,
$0\le b_j\le B_j$ per tutti i $j$ e
$\sum_j e_j^2/V_j\le c_R^2$.
Ogni somma relativa a una coppia consecutiva è limitata dalla somma totale.
Cauchy--Schwarz dà perciò $M_r-tM_{r-1}\ge F_{r,R}(t)$ simultaneamente
per tutti gli $r$ e tutti i $t\ge0$. Il ragionamento del Teorema
\ref{thm:pair} vale ora su un solo evento e protegge anche il massimo in $r$.
\end{proof}
\begin{remark}[Che cosa è e non è adattivo]
La soglia $t$ e l'ordine $r\in\{2,\ldots,R\}$ possono essere scelti
dopo aver osservato i dati. Il massimo $R$ e il livello $\delta$ sono
fissati prima del campione. Un ulteriore confronto tra diversi $R$
richiederebbe un altro controllo simultaneo. Non è giustificato
massimizzare soglie a ordine fissato senza la calibrazione comune.
\end{remark}

\section{Corrispondenza con l'implementazione e lettura dei risultati}
La procedura simultanea usa soltanto $n$, $\delta$, $R$ e
$K_1,\ldots,K_{2R}$. Per $R=6$ occorrono i primi dodici conteggi dello
spettro. Se $K_r=0$, allora $T_r=0$, $F_{r,R}(0)<0$ e la soglia
di quell'ordine è zero. Un conteggio $K_{2r}=0$ non rende nulla la
varianza superiore: $C_{2r}>0$.
Gli estremi numerici vengono arrotondati verso il lato ammissibile
della radice. Il vecchio bound e $\Lambda_6$ sono riportati separatamente;
il loro massimo non è automaticamente valido allo stesso livello.
La copertura Monte Carlo di un endpoint $L$ è
\[
 B_{\mathrm{rep}}^{-1}\sum_{b=1}^{B_{\mathrm{rep}}}
 \ind\{L_b\le M_{\max,b}\},
\]
non la frequenza di confronto con l'oracle deterministico.
Non si asserisce che $\Lambda_6$ domini il vecchio bound: il vantaggio
geometrico può essere superato dalla penalità probabilistica di Markov.

\begin{thebibliography}{9}
\bibitem{ayed2019} F. Ayed, M. Battiston, F. Camerlenghi, S. Favaro (2019).
A Good--Turing estimator for feature allocation models.
\emph{Electronic Journal of Statistics} 13, 3775--3804.
\end{thebibliography}
\end{document}
